\cotefile{19}
\scanpage{2}
Résolutions standarts et foncteurs adjoints

Tapis Cartier

\margin{$9_{\mathrm p}+14_{\mathrm p}$}
% archive sheet number: 1

\scanpage{4}
\uline{Tapis Cartier}

1)
\diagram{\begin{tikzcd}[ampersand replacement=\&]
 A \arrow[r,shift left,"f"] \& B \arrow[l,shift left,"g"]
\end{tikzcd}}
$g$ adjoint : \deleted{\unsure{gauche}}droite de $f$, d'où morphismes d'adjonction
\[
 \mathrm{id}_A\xrightarrow{\beta}gf,\qquad fg\mathrel{\overset{\alpha}{\unsure{-|\to}}}\mathrm{id}_B \tag*{(1.1)}
\]
On pose
\[
 \varphi=fg:\ B\mathrel{\unsure{-|\to}} B \tag*{(1.2)}
\]
On a donc
\[
 \boxed{\varphi\xrightarrow{\alpha}\mathrm{id}_B,\quad\varphi\xrightarrow{\lambda}\varphi^2} \tag*{(1.3)}
\]
où
\[
 \lambda\deleted{\unsure{=}}:fg\longrightarrow fgfg\qquad \lambda=f*\beta*g \tag*{(1.4)}
\]
\deleted{Alors} $\varphi$ satisfait \uline{conditions de Godement}

\diagram{\begin{tikzcd}[ampersand replacement=\&,column sep=small]
 \text{(1.5) (\uline{unitarité})} \& \varphi \arrow[r,"\lambda"] \arrow[rr,bend right=30,"\mathrm{id}_{\varphi}"'] \&
 \varphi^2 \arrow[r,shift left,"\varphi*\alpha\deleted{\unsure{*\varphi}}"] \arrow[r,shift right,"\alpha*\varphi"'] \& \varphi
\end{tikzcd}}
commutatif : $(\varphi*\alpha)\circ\lambda=\mathrm{id}_{\varphi}$

\diagram{\begin{tikzcd}[ampersand replacement=\&,column sep=small]
 \text{(1.6) (\uline{associativité})} \& \varphi\arrow[r,"\lambda"] \& \varphi^2
 \arrow[r,shift left,"\varphi*\lambda"] \arrow[r,shift right,"\lambda*\varphi"'] \& \varphi^3
\end{tikzcd}}
commutatif : $(\varphi*\lambda)\circ\lambda=(\lambda*\varphi)\circ\lambda$

On définit la catégorie
\[
 K=K(\varphi,\alpha,\lambda) \tag*{(1.7)}
\]
comme formée des couples
\[
 (X,u)\qquad X\in\mathrm{Ob}\,B,\quad \boxed{u:X\mathrel{\unsure{-|\to}}\varphi(X)} \tag*{(1.8)}
\]
satisfaisant les conditions de commutation
\diagram{\resizebox{\linewidth}{!}{\(\left\{\begin{tikzcd}[ampersand replacement=\&,column sep=small,row sep=large]
 \text{(\deleted{\unsure{1.9}}1.10)} \& {}\arrow[d,phantom,"{\left\{\vphantom{\rule{0pt}{2cm}}\right.}" description] \& X\arrow[r,"u"] \& \varphi(\deleted{\unsure{Y}}X)
 \arrow[r,shift left,"\varphi(u)"] \arrow[r,shift right,"\lambda(X)"'] \& \varphi^2(X)
 \& \lambda\circ u=\varphi(u)\circ u \quad\text{``associativité''}\\
 \text{(\deleted{\unsure{1.10}}1.9)} \& {} \& X\arrow[r,"u"] \& \varphi(X)\arrow[r,"\alpha(X)"] \& X
 \& \alpha(X)\circ u=\mathrm{id}_X\quad\text{``unitarité''}
\end{tikzcd}\right.\)}}
On obtient un foncteur naturel
\[
 \boxed{h:A\longrightarrow K(\varphi,\alpha,\lambda)} \tag*{(1.11)}
\]
par
\[
 Y\longmapsto \bigl(\deleted{\unsure{f}}f(Y),\quad f(Y)\xrightarrow{f*\beta}\varphi f(Y)=fgf(Y)\bigr).
\]
\margin{vérifier compatibilité (1.9)}
% archive sheet number: 3

\scanpage{5}
\uline{Théorème} 1.12 (Cartier) \deleted{Supposons que} \inserted{a) Si} dans $A$ les $\mathrm{Ker}$ de doubles flèches existent,$^{*}$ et \deleted{que} \inserted{si} $f$ y commute $(*)$, alors 1.11 est essentiellement surjectif,
% slanted
\margin{$*$ il suffit de regarder les doubles flèches $u,v$ telles que $\mathrm{Ker}(f(u),f(v))$ \uline{existe}}
\deleted{\unsure{et}} \inserted{b)} Pour que \uline{$h$} soit une équivalence [i.e. soit \deleted{de la forme} pleinement fidèle], il f. et s. que $f$ soit \uline{conservatif}, et que pour toute double flèche \inserted{$(u,v)$} de $A$ telle que $\mathrm{Ker}(f(u),f(v))$ existe, on ait que $\mathrm{Ker}(u,v)$ existe \uline{et} $f$ y commute.
% slanted
\margin{(condition \unsure{c|C})}
% slanted
\margin{améliorer a), en explicitant que $h$ est un \uline{passage au quotient}}

2) Partons d'une catégorie $B$, avec un foncteur
\[
 \varphi:B\mathrel{\unsure{-|\to}}B, \tag*{(2.1)}
\]
\deleted{satisfaisant les conditions (1.5),} et des liaisons \unsure{(1.2)}
\deleted{(2.1}
\[
 \varphi\mathrel{\overset{\alpha}{\unsure{-|\to}}}\mathrm{id}_B,\qquad\varphi\xrightarrow{\lambda}\varphi^2 \tag*{(2.2)}
\]
satisfaisant les conditions (1.5), (1.6). On peut alors (\(=\) sans ces restrictions (1.5) et (1.6) !) définir une catégorie
\[
 A=K(\varphi,\alpha,\lambda) \tag*{(2.3)}
\]
comme ci-dessus. Soit
\[
 A\xrightarrow{f}B \tag*{(2.4)}
\]
le foncteur \deleted{``de \unsure{réduction|oubli}'' \unsure{en}} canonique $(X,u)\longmapsto X$. On \deleted{\uline{Question} Quand $f$ admet-il un adjoint à droite $g$, de telle façon que \unsure{l'on ait} $\varphi\simeq fg$} va définir un \deleted{\unsure{objet}}foncteur
\[
 g:B\longrightarrow A \tag*{(2.5)}
\]
par
\[
 g(X)=\bigl(\varphi(X),\quad\varphi(X)\xrightarrow{\lambda(X)}\varphi^{2}(X)\bigr) \tag*{(2.6)}
\]
[Les conditions \inserted{\textcircled{\scriptsize 1.5}} \textcircled{\scriptsize (1.6)} assurent qu'on tombe bien dans $K(\varphi,\alpha,\lambda)$ !].
\deleted{Je dis qu'on a une relation d'adjonction,}
\[
 \deleted{\mathrm{Hom}_A\bigl((X,u),g(Y)\bigr)\simeq\mathrm{Hom}_B\bigl(\underset{\shortparallel\,X}{f(X,u)},\unsure{Y|X}\bigr)} \tag*{\deleted{(2.7)}}
\]
et on a évidemment
\[
 fg=\varphi \tag*{(2.7)}
\]
% archive sheet number: 4

\scanpage{6}
Considérons alors
\[
 fg\xrightarrow{\alpha}\mathrm{id} \tag*{(2.8)}
\]
ce qui fournit un homomorphisme d'adjonction
\[
 \begin{aligned}
 &\inserted{v\longmapsto\alpha\deleted{\unsure{(Y)}}\circ f(v):}\\
 &\deleted{\unsure{\mathrm{Hom}(fX,Y)\longleftarrow}}\alpha':\mathrm{Hom}_A(X,g(Y))\longrightarrow\mathrm{Hom}_B(f(X),Y)
 \end{aligned} \tag*{(2.9)}
\]
\deleted{\unsure{par}} On va définir aussi
\[
 \deleted{v\longmapsto\alpha\circ f(\unsure{u|v})} \tag*{\deleted{(2.9')}}
\]
\[
 \deleted{\unsure{g}}\beta:\mathrm{id}_K\longrightarrow gf \tag*{(2.\deleted{\unsure{9}}10)}
\]
\unsure{par|en} \unsure{prenant|posant}, pour tt $Y=(X,u)$ dans $K$,
\diagram{\begin{tikzcd}[ampersand replacement=\&]
 X\arrow[r,"u"]\arrow[d,"u"'] \& \varphi(X)\arrow[d,dashed,"\varphi(u)"] \& \unsure{\leftarrow}Y,\\
 \varphi(X)\arrow[r,"\lambda(X)"'] \& \varphi^2(X) \& g\deleted{\unsure{g}}f(Y)
\end{tikzcd}}
la commutativité étant assurée par la \uuline{condition} \textcircled{\scriptsize (1.\deleted{\unsure{9}}10)} sur les objets de $K$. On a ainsi un hom. d'adjonction en l'autre sens
\[
 \beta':\mathrm{Hom}_B(fX,Y)\longrightarrow\mathrm{Hom}_A(X,g(Y))
 \qquad v\longmapsto g(v)\circ\beta(X) \tag*{(2.11)}
\]
On vérifie que ces hom. sont \uline{inverses l'un de l'autre}. En effet, $\beta'\alpha'=\mathrm{id}$ résulte facilement de \textcircled{\scriptsize 1.5}, et $\alpha'\beta=\mathrm{id}$ \inserted{résulte de \textcircled{\scriptsize (1.\deleted{\unsure{10}}9)}.}
\deleted{équivaut à ceci : Pour \unsure{tout} objet $(X,u)$ de $A$, le composé}
\diagram{\begin{tikzcd}[ampersand replacement=\&,every arrow/.append style={dash pattern=on 2pt off 0pt}]
 \deleted{X}\arrow[r,"{\deleted{u}}"] \& \deleted{\varphi(X)}\arrow[r,"{\deleted{\alpha(X)}}"] \& \deleted{X}
\end{tikzcd}}
\deleted{est \uline{l'identité}. Or \unsure{on a} (1.9) $u=\mathrm{id}_{\varphi(X)}u
 \overset{1.5}{=}\varphi(\alpha(X))\lambda(X)u\overset{1.9}{=}\varphi(\alpha(X))\varphi(u)u
 =\varphi(\alpha(X)u)$ et (1.5) $\varphi(u)\circ u=\lambda(X)\circ u:X\longrightarrow\varphi^2(X)$
 d'où, en composant à gauche avec $\lambda(X)$
 $\underbrace{\lambda(X)\varphi(u)}\circ u=\underbrace{\lambda(X)\lambda(X)}u$}

Donc $f$ et $g$ sont adjoints l'un de l'autre, et \uline{redonnent bien sûr} \deleted{\unsure{par le}} $(\varphi,\alpha,\lambda)$ \uline{par le procédé du} N° 1.

Le foncteur $f$ est conservatif et satisfait \unsure{?|c|e} \ldots
% archive sheet number: 5

\scanpage{7}
\uline{Conclusion} La donnée (pour $B$ fixé) d'une catégorie $A$, et d'un couple de foncteurs adjoints $f,g$,
\diagram{\begin{tikzcd}[ampersand replacement=\&] A\arrow[r,shift left,"f"] \& B\arrow[l,shift left,"g"]\end{tikzcd}}
avec $f$ \uline{conservatif} et \deleted{tel que, \unsure{pour les}} \deleted{commutant aux noyaux de doubles \unsure{doubles} flèches} satisfaisant la condition C de 1.12 est équivalente essentiellement à la donnée d'un foncteur $\varphi$ dans $B$, avec les données de Godement $\alpha,\lambda$.

\deleted{De plus, $f$ est pleinement fidèle ssi $\alpha$ est un isomorphisme,}

\uline{Corollaire} $f$ pleinement fidèle \deleted{\unsure{\Longleftrightarrow}}\quad
$\alpha*\varphi\text{ isom}\Longleftrightarrow\lambda\text{ isom}$ [\unsure{alrs|NB} alors $\lambda$ est connu quand on connaît $\alpha$]

\uline{Compléments sur propriétés d'exactitude}
% slanted
\margin{\deleted{Si de plus les \unsure{noyaux}} \uline{Pour} les couples $(f,g)$ \deleted{\unsure{de foncteurs}} adjoints, où \unsure{$\exists$ ds $B$} les Ker \unsure{resp.}, $f$ conservatif et commute aux Ker}

a) Pour tout type de $\varprojlim$ qui sont représentables dans $B$, \deleted{et auxquelles $\varphi$ commute,} les $\varprojlim$ de ce type \deleted{existent} sont repr. dans $A$, et le foncteur $f$ y commute, \inserted{si et seulement si $\varphi$ y commute.}

b) Pour tout type de $\varinjlim$ qui est représentable dans $B$, \unsure{ce} \unsure{\textquotedbl} type est représentable dans $A$, et $f$ y commute (il commute \textquotedbl à \uline{toutes} les $\underline{\varinjlim}$ !)

c) Pour tt propriété d'exactitude de $B$, \deleted{\unsure{faisant|faisant intervenir}} du type commutation de $\varprojlim_J$ \deleted{de type $J$} à $\varinjlim_I$, avec des \deleted{\unsure{typ.}} $\lim$ en question repr. dans $B$, et $\varphi$ comm. aux $\varprojlim_J$, -- la \textquotedbl propriété d'exactitude est vraie dans $A$. \uline{Ex.} Les propriétés d'exactitude a) b) c)

\deleted{d) Supposons que $B$ soit une \unsure{filtration|catégorie} \inserted{du topos} \unsure{continue} $(B_k)$,}
\deleted{\unsure{d)}}
\uline{Théorème} Supposons $B$ un topos. Alors $A$ est un topos et $f$ de la forme $f_0^*$ pour un morphisme de topos $B\longrightarrow A$ (nec. conservatif) sss $\varphi$ \deleted{commute aux \unsure{$\varprojlim$}} \inserted{est exact à gauche \uline{et accessible}}. Pour que $f_0$ soit essentiel, il f. et s. que $\varphi$ commute aux $\varprojlim$ quelconques i.e. aux produits infinis, (ou encore que $\varphi$ ait un adjoint à gauche $\psi$).
% archive sheet number: 6

\scanpage{8}
3. Supposons
\[
 \boxed{B=\prod_{i\in I}B_i}. \tag*{(3.1)}
\]
Donc $f$ est donné par une famille de
\[
 f_i:A\longrightarrow B_i \tag*{(3.2)}
\]
et $g$ est connu quand on connaît les adjoints
\[
 g_i:B_i\longrightarrow A, \tag*{(3.3)}
\]
par
\[
 g\bigl((X_i)_{i\in I}\bigr)=\prod_{i\in I}g_i(X_i) \tag*{(3.4)}
\]
Le composé $\varphi=fg$ est alors un \unsure{hom|hom.}
\[
 \varphi:\prod B_i\longrightarrow\prod B_i
\]
et est connu quand on connaît les
\[
 \mathrm{pr}_{j}\,\varphi:\prod_i B_i\mathrel{\unsure{-|\to}}B_j
\]
On trouve
\[
 \mathrm{pr}_j\,\varphi((X_i))=f_j\left(\prod_{i\in I}g_i(X_i)\right). \tag*{(3.5)}
\]
Soit
\[
 \varphi_{ji}\overset{\text{déf}}{=}f_jg_i\deleted{\unsure{g_j f_i}}:\ B_i\longrightarrow B_j \tag*{(3.6)}
\]
et supposons que les $f_j$ commutent aux produits de type $I$ (p. ex. $f_j$ exacts à g. et $I$ fini). On a alors
\[
 \mathrm{pr}_j\,\varphi((X_i))\simeq\prod_{i\in I}\varphi_{ji}(X_i) \tag*{(3.6)}
\]
et les données de $\alpha,\lambda$ reviennent à celles de
\[
 \prod_i\varphi_{ki}(X_i)\xrightarrow{\bar\alpha_{\deleted{\unsure{j}}k}}X_{\deleted{\unsure{j}}k} \tag*{(3.7)}
\]
\[
 \prod_i\varphi_{\deleted{\unsure{j}}ki}(X_i)\xrightarrow{\bar\lambda_k}
 \prod_{\deleted{\unsure{j}}i,j}\varphi_{kj}\varphi_{ji}(X_i) \tag*{(3.8)}
\]
i.e.
\[
 \prod_\alpha\varphi_{k\alpha}(X_\alpha)\xrightarrow{\bar\lambda_{k,j,i}}\varphi_{kj}\varphi_{ji}(X_i)
\]
% archive sheet number: 7

\scanpage{9}
qui s'explicitent, en termes des $\varphi_i=f_i g_i$, et des $\alpha_i,\lambda_i$ correspondants, ainsi
\[
 \left\{\begin{aligned}
 \bar\alpha_k&=\alpha_k(X_k)\circ\mathrm{pr}_k\\
 \bar\lambda_{k,j,i}&=\lambda_{kji}(X_i)\circ\mathrm{pr}_i
 \end{aligned}\right. \tag*{(3.9)}
\]
\diagram{\((3.10)\quad\boxed{\begin{tikzcd}[ampersand replacement=\&,column sep=1.3em,row sep=.25em]
 \lambda_{kji}: \& \varphi_{ki}\arrow[r] \& \varphi_{kj}\circ\varphi_{ji}\\
 \& \shortparallel \& \shortparallel\\
 \& f_k g_i \& f_k\underbrace{g_j f_j}g_i\\
 \& \& \mathrm{id}\arrow[u]
 \end{tikzcd}}\)}
Ces $(\varphi_i,\lambda_{ijk})$ sont reliés par des conditions exprimant (1.5) (1.6), que nous n'explicitons pas. \inserted{Pour \deleted{\unsure{k}}$i=j=k$, $\lambda_{kji}$ n'est autre que $\lambda_i$.} Le cas le plus intéressant est celui où les \deleted{\unsure{$\varphi_i$|$\varphi$}} $\alpha_i$ sont des isom, (i.e. les $g_i$ \unsure{pl|plt} fidèles).
% slanted
\margin{notons que si deux des indices $i,j,k$ coïncident, alors $\lambda_{kji}$ est trivialement donné en termes des $\lambda_i$ --}

Dans ce cas, \deleted{il est} on peut supposer $\varphi_i=\mathrm{id}_{B_i}$, $\alpha_i=\mathrm{id}_{\mathrm{id}_{B_i}}$, $\lambda_i=\mathrm{id}$. On constate (!) que \inserted{dans} le cas où deux \inserted{consécutifs} des indices $i,j,k$ sont égaux, i.e. $i=j$ ou $j=k$, de sorte que l'un des facteurs de $\varphi_{kj}\varphi_{ji}$ est l'identité, on trouve \inserted{$\lambda_{kji}$} l'\uline{identité}.

Il reste donc à examiner les cas
\[
 \left\{\begin{aligned}
 \text{a)}\ &\deleted{\unsure{i\ne j,\ j\ne k,\ i\ne{}}}\quad i,j,k\ \text{tous différents}\\
 \text{b)}\ &i=k\ne j\qquad\mathrm{id}_{B_i}\longrightarrow\varphi_{ij}\circ\varphi_{ji}\qquad(i\ne j)
 \end{aligned}\right. \tag*{(3.11)}
\]
\deleted{Si \unsure{nous}}

\centerline{\rule{.72cm}{.3pt}\quad\rule{.72cm}{.3pt}\quad\rule{.72cm}{.3pt}\quad\rule{.9cm}{.3pt}}
% archive sheet number: 8

\scanpage{10}
\deleted{avec des conditions (1.5) (1.6) qu'on n'explicite pas.}

Prenons le cas (3.12) $B=B'\times B''$, $f=(f',f'')$, $g(X',X'')=g'(X')\times g''(X'')$, \deleted{alors les données sont, si on pose concernant $\alpha$} supposons \unsure{que|de plus} $f'g'\simeq\mathrm{id}_{B^{\unsure{''|'}}}$, \deleted{\unsure{$f'g'$}} $f''g''\simeq\mathrm{id}_{B''}$ (i.e. $g',g''$ pl. fidèles), donc $\varphi$ est connu quand on connaît
\[
 \boxed{\varphi'':B'\longrightarrow B'',\qquad \varphi':B''\mathrel{\unsure{-|\to}}B'} \tag*{(3.13)}
\]
\deleted{et la donnée des $\alpha$ \unsure{revient à} $\lambda_{kji}$ revient à celle des}
\[
 \deleted{X'\times\varphi'(\deleted{\unsure{B}}X'')\longrightarrow X'} \tag*{\deleted{(3.10)}}
\]
\[
 \deleted{\varphi''(X')\times X''\longrightarrow X''} \tag*{\deleted{(3.11)}}
\]
\[
 \deleted{\left\{\begin{aligned}
 X'&\longrightarrow\varphi'\varphi''(X')\qquad X''\mathrel{\unsure{-|\to}}\varphi''\varphi'(X'')\\
 &X
 \end{aligned}\right.}
\]
\deleted{Il y a huit $\lambda_{kji}$ ($2^3$), dont deux triviaux ($i=j=k$), donc reste 6 non triviaux $\lambda_{1,1,2}:$} Le cas a) ci-dessus ne peut se présenter, et il reste donc des données
\[
 \boxed{\left\{\begin{aligned}
 \mathrm{id}_{B'}&\xrightarrow{\lambda'}\varphi'\varphi''\\
 \mathrm{id}_{B''}&\xrightarrow{\lambda''}\varphi''\varphi'
 \end{aligned}\right.} \tag*{(3.14)}
\]
Il faudrait encore \deleted{vérifier} expliciter dans ce cas les conditions (1.5) et (1.6) --

\diagram{\resizebox{\linewidth}{!}{\begin{tikzcd}[ampersand replacement=\&,column sep=small]
 \varphi(X',X'')\arrow[r,"\lambda"] \& \varphi^2(X',X'')\arrow[r,shift left,"\text{sont égaux ?}"]\arrow[r,shift right] \& \varphi(X',X'')\\[-8pt]
 \shortparallel \& \shortparallel \& \shortparallel\\[-8pt]
 (X'\times\varphi'(X''),X''\times\varphi''(X'))\arrow[r] \&
 \begin{gathered}\bigl(X'\times\varphi'(X'')\times\varphi'(X'')\times\varphi'\varphi''X',\\[-2pt]
 X''\times\varphi''(X')\times\varphi''(X')\times\varphi''\varphi'(X'')\bigr)\end{gathered} \&
 \begin{gathered}(X'\times\varphi'(X''),\\ X''\times\varphi''(X'))\end{gathered}
\end{tikzcd}}}
\margin{cf. 1.5 semble O.K.}
\deleted{Par contre, il semble que 1.6 impose des restrictions à $\lambda',\lambda''$}
Il \deleted{\unsure{ne}} semble \deleted{pas} qu'il n'y ait \deleted{pas} les conditions \ldots
\[
 \boxed{\begin{aligned}
 \varphi''*\lambda'&=\deleted{\unsure{\varphi}}\lambda''*\varphi'':\varphi''\longrightarrow\varphi''\varphi'\varphi''\,;\\
 \lambda'*\varphi'&=\varphi'*\lambda'':\varphi'\deleted{\unsure{\varphi}}\longrightarrow\varphi'\varphi''\varphi'
 \end{aligned}}
\]
Explicitons la catégorie $A$ en termes de $\varphi',\varphi'',\lambda',\lambda''$

Elle est formée des couples $(X',\unsure{\alpha''|X''},u',u'')$
\[
 X'\xrightarrow{u'}\varphi'(X''),\qquad X''\xrightarrow{u''}\varphi''(X') \tag*{(3.15)}
\]
\deleted{$X'\xrightarrow{\unsure{u'}}$} définissant
% archive sheet number: 9

\scanpage{11}
\deleted{soumis aux conditions}

(définissant \deleted{\unsure{$X'\longleftarrow(\mathrm{id}_{X'}$}} $X'\xrightarrow{(\mathrm{id}_{X'},u')}X'\times\varphi'(X'')$, $X''\xrightarrow{(\mathrm{id}_{X''},u'')} \deleted{\unsure{X'}}X''\times\varphi''(X')$\deleted{\unsure{e}}, ce qui exprime déjà (1.9)) soumis à la condition de commutativité de (1.10)
\diagram{\resizebox{\linewidth}{!}{\begin{tikzcd}[ampersand replacement=\&,column sep=large]
 (X',X'')\arrow[r,"*","{\text{via }u',u''}"'] \&
 (X'\times\varphi'(X''),X''\times\varphi''(X'))\arrow[r,shift left,"{\varphi\,(u)}"]\arrow[r,shift right,"{\substack{\text{can}\\\text{via }\lambda',\lambda''}}"'] \&
 \begin{gathered}\bigl(X'\times\varphi'(X'')\times\varphi'(X'')\times\varphi'\varphi''(X'),\\
 X''\times\varphi''(X')\times\varphi''(\unsure{X'|X''})\times\varphi''\varphi'(X'')\bigr)\end{gathered}
\end{tikzcd}}}
\[
 \deleted{\unsure{X'}\quad X'\xrightarrow{\mathrm{id},u'}X'\times\varphi'(X'')\longrightarrow}
\]
il semble qu'on trouve la commutativité des diagrammes :
\diagram{\(\boxed{\text{(3.16)}}\left\{\begin{tikzcd}[ampersand replacement=\&,row sep=huge]
 X'\arrow[r,"u'"]\arrow[rr,bend right=20,"{\lambda'(X')}"'] \& \varphi'(X'')\arrow[r,"{\varphi'(u'')}"] \& \varphi'\varphi''(X')\\
 X''\arrow[r,"u''"]\arrow[rr,bend right=20,"{\lambda''(X'')}"'] \& \varphi''(X')\arrow[r,"{\varphi''(u')}"] \& \varphi''\varphi'(X'')
\end{tikzcd}\right.\)}

\uline{Exemple} Si $\varphi''\varphi'$ et $\varphi'\varphi''$ sont des foncteurs constants de valeurs un objet final de $B'$ resp. $B''$, les conditions 3.16 sont vides, et on trouve simplement --.

\[
 \resizebox{.91\linewidth}{!}{\(\displaystyle\left\{\begin{aligned}
 &\left(X',X'',X'\xrightarrow{u'}\varphi'(X''),X''\xrightarrow{u''}\varphi''(X')\right)
 \quad \begin{aligned}&\overset{f'}{\longmapsto}X'\\[-2pt]&\xrightarrow{f''}X''\end{aligned}\\
 &g'(X')=\left(X',\varphi''(X'),X'\xrightarrow{\lambda'(X')}\varphi'\varphi''(X'),\varphi''(X')\xrightarrow{\mathrm{id}}\varphi''(X')\right)\\
 &g''(X'')=\left(\varphi'(X''),X'',\varphi'(X'')\xrightarrow{\mathrm{id}}\varphi'(X''),\deleted{\unsure{\varphi}}X''\xrightarrow{\lambda''(X'')}\varphi''\varphi'(X'')\right)
 \end{aligned}\right.\)} \tag*{3.17}
\]
\deleted{\unsure{$g'f'$}}
\[
 \left\{\begin{aligned}
 \text{Image ess. de }g' &: \quad u''\ \text{\uline{un isom}}\\
 \text{Image ess. de }g'' &: \quad u'\ \text{\uline{un isom}}
 \end{aligned}\right. \tag*{3.18}
\]
\margin{$\unsure{\perp|1}$}
% archive sheet number: 10

\scanpage{12}
(3.20)
\diagram{\resizebox{\linewidth}{!}{\begin{tikzcd}[ampersand replacement=\&,column sep=small,row sep=huge]
 X\arrow[d,"{\beta'}"'] \& = \& (X'\arrow[d,"{\textcircled{\scriptsize id}}"'] \& ,X''\arrow[d,"{\textcircled{\scriptsize $u''$}}"] \& ,X'\arrow[r,"u'"]\arrow[d,no head,"\mathrm{id}"] \& \varphi'(X'')\arrow[d,no head,"{\varphi'(u'')}"] \& ,X''\arrow[r,"{\deleted{\unsure{\varphi}}u''}"]\arrow[d,"u''"] \& \varphi''(X'))\arrow[d,no head,"\mathrm{id}"]\\
 g'f'(X) \& = \& (X' \& ,\varphi''(X') \& ,X'\arrow[r,"{\lambda'(X')}"] \& \varphi'\varphi''(X') \& ,\varphi''(X')\arrow[r,"\mathrm{id}"] \& \varphi''(X'))
\end{tikzcd}}}
\margin{adjonction}
et $\beta''$ analogue

[\uline{Question} Quand $\beta'$ est-il un \uline{épimorphisme} pour tout $X$ ? Cela \deleted{\unsure{implique|signifie}} \inserted{signifie}, moyennant des conditions d'exactitude : \unsure{droit|disons|droite} \unsure{standart|standard} sur $\varphi',\varphi''$, que $u'':X''\longrightarrow\varphi''(X')$ est tjrs un épimorphisme. ]

L'intersection des images essentielles de $g'$ et de $g''$ est formée des $(X',X'',u',u'')$ pour lesquels \deleted{\unsure{$g'$}}$u'$ \uline{et} \deleted{\unsure{$g''$}}$u''$ sont des isom. \deleted{Cette catégorie \unsure{i.}} Cette catégorie équivaut à l'image inverse par $g'$ de $g''_{\mathrm{ess}}(B'')$, qui est la sous-catégorie \inserted{str} pleine de $B'$ formée des $X'$ tels que $\lambda'(X')$ soit un isom. On trouve donc

(3.21) Conditions équivalentes :

(i) $g'_{\mathrm{ess}}(B')\wedge g''_{\mathrm{ess}}(B'')\simeq\text{catégorie finale}$

(ii) La sous-catégorie pleine de $B'$ formée des $X'$ tels que $\lambda'(X')$ soit un isom, est équivalente à la \deleted{sous}-catégorie \unsure{finale|final}

(ii bis) Idem pour $\lambda''$
% archive sheet number: 11

\scanpage{14}
\uline{\unsure{Résolutions|Résolution} standard et théories relatives.}

\diagram{\begin{tikzcd}[ampersand replacement=\&]
 C'\arrow[r,shift left,"G"] \& C\arrow[l,shift left,"F"]
\end{tikzcd}}
% slanted
\margin{\uline{\unsure{Dont|Droit|Donné}}}
% slanted
\margin{\unsure{Godement|G-préfaisc.|G-faisc.}}
\[
 \mathrm{Hom}(F(X),X')\xrightarrow{\sim}\mathrm{Hom}(X,G(X'))
 \qquad F\quad G
\]
Posons
\[
 E=GF:C\longrightarrow C
\]
d'où
\[
 \varphi:\mathrm{id}_C=E^0\longrightarrow E,\qquad \chi:E^2\longrightarrow E
\]
satisfaisant les relations de Godement, d'où un objet \inserted{\textcircled{\scriptsize co}}\uline{simplicial} \deleted{dans} $\widehat E$ dans \uline{Hom} $(C,C)$ \deleted{\unsure{$\varphi$ et $\chi$ où}}.
\[
 \widehat E^{\,n}=E^{n+1}\ \deleted{\unsure{n}}\ \ldots\ldots\quad
 \text{et \uline{augmentation}}\quad X\longrightarrow\widehat E(X)
\]
(On devrait trouver une interprétation naturelle par théorie de la descente). On pose de \textquotedbl
\[
 E'=FG:C\longrightarrow C
\]
d'où
\[
 \varphi':E'\longrightarrow\mathrm{id}_{C'}=E'^{\,0},\qquad
 \chi':E'\mathrel{\unsure{-|\to}}E'^{\,2}
\]
satisfaisant aux relations de Godement, d'où un objet \uline{simplicial} $\widehat E'$ de \uline{Hom} $(C',C')$ :
\[
 \widehat E'^{\deleted{\unsure{n}}}_{n}=E'^{\,n+1}\qquad\text{---}
\]
Notons les \deleted{\unsure{isomorphismes} \unsure{fonctoriels|simpliciaux}} \inserted{\unsure{résultats suivants}}
\[
 \begin{aligned}
 \text{(i)}\quad &G(X')\mathrel{\substack{\sim\\[-.7ex]\sim\\[-.7ex]\sim}}\widehat E(G(X'))\ \deleted{\unsure{\simeq0}}
       &&\text{pour tt }X'\in\mathrm{ob}\,C'\\
 \text{(ii)}\quad &F(X)\approxeq F(\widehat E(X))\ \deleted{\unsure{\simeq0}}
       &&\text{pour tt }X\in\mathrm{ob}\,C
 \end{aligned}
\]
La deuxième relation est vraie plus gén. si on a un foncteur $T:C\longrightarrow D$ et un hom \unsure{fonctoriel}
\[
 h:T\circ E\longrightarrow T
\]
tel que
\[
 h\circ(T*\varphi)=\mathrm{id}_T
\]
alors
\[
 T\widehat E(X)\approxeq 0\qquad\text{pour tt }X\in\mathrm{ob}\,C
\]
[cf Godement]. La formule (i) se vérifie sans doute de \textquotedbl par des formules explicites.

\uline{Application} Supposons $C,C'$ abéliennes, $F,\unsure{G|C}$ additifs, \deleted{\illegible{2}} \deleted{\illegible{1}}. a) Si $F$ est \uline{exact} et \uline{fidèle}, alors $\widehat E(X)$ est une \uline{résolution} de $X$ pour tt $X$.

b) Si $\Gamma$ est un foncteur
% archive sheet number: 13

\scanpage{15}
\uline{Cohomologies relatives} (un cas spécial)

1. \uline{Caractérisation des foncteurs} $A\longrightarrow\mathrm{Hom}(A,B)$.

Soit $C$ une catégorie. Pour tout $B\in C$, soit $h_B$ le foncteur de $C$ dans les ensembles tel que
\[
 h_B(A)=\mathrm{Hom}(A,B)
\]
$h_B$ est un foncteur \unsure{.l|,!|,} contravariant en $B$, à valeurs dans la catégorie des foncteurs $C\longrightarrow\unsure{\mathcal E|E}$, et on a
\[
 \alpha:\mathrm{Hom}_C(B,B')\xrightarrow{\sim}\mathrm{Hom}(h_B,h_{B'}) \tag*{(*)}
\]
[En effet, on définit la flèche inverse $\beta$, en associant à \deleted{\unsure{tout}} $u\in\mathrm{Hom}(h_B,h_{B'})$ \deleted{\unsure{un}} l'élément $\beta(u)$ de $\mathrm{Hom}_C(B,B')$ image par $u_B$ de l'élément unité de $h_B(B)=\mathrm{Hom}_{\unsure{C|a}}(B,B)$ dans $h_B(B')=\mathrm{Hom}_C(B,B')$].

\deleted{Si dans $C$ il y a \unsure{assez de sommes}} Si $C$ est additive, alors $h_B$ est additif \inserted{\uline{exact à g.}}, le foncteur $B\longrightarrow h_B$ est additif (exact \uline{à gauche}) et l'isomorphisme $(*)$ est valable en se limitant aux hom. \uline{additifs}. Ainsi $B\longrightarrow h_B$ \deleted{\unsure{est}} permet d'``identifier'' $C$ à une sous catégorie de la catégorie $\mathrm{Fonc.add.}(C^o,\unsure{\mathscr G|\mathfrak G|\mathcal Y})$.

Supposons que $C$ ait un générateur $U$, et \deleted{que} satisfasse à AB 3, 4, 5, . \deleted{\unsure{Comme|Le foncteur}} \inserted{\unsure{Tout|Tjs est}} \deleted{\unsure{alors}} $h_B$ transforme limites inductives en limites projectives, et \deleted{\unsure{par suite} \unsure{aussi}} \inserted{est exact} à gauche. On va donner une réciproque. Tout d'abord, si $h,h'$ sont deux foncteurs \inserted{contr.} $C^{\deleted{\unsure{o}}}\longrightarrow\mathscr G$ ayant
% archive sheet number: 14

\scanpage{16}
ces propriétés, on montre que, posant $M(h)=h(U)$, on \deleted{\unsure{est}} a
\[
 \alpha':\mathrm{Hom}(h,h')\xrightarrow{\sim}\mathrm{Hom}_{R^o}(\deleted{\unsure{h}}M(h),M(h'))
\]
où $R=\mathrm{Hom}_C(U,U)$, $R^o=$ anneau inverse de $R$.

Pour le voir, on définit un hom
\[
 \beta':\mathrm{Hom}_{R^o}(M(h),M(h'))\longrightarrow\mathrm{Hom}(h,h')
\]
Détails laissés au lecteur. Donc le foncteur $T':h\longrightarrow M(h)$ identifie la catégorie des foncteurs \deleted{Posons donc, pour \unsure{les}} envisagés à une sous-catégorie de $\unsure{\mathscr G|\mathfrak G|\mathcal Y}^{R^o}$ ($R$-mod à droite). On va définir un ``foncteur inverse'' $M\longrightarrow h_M$ de $\mathscr G^{R^o}$ dans la catégorie de foncteurs envisagée. On va définir directement le foncteur $T'':\deleted{\unsure{C}}\mathscr G^{R^o}\longrightarrow C$, ce sera simplement
\[
 T''(M)=M\otimes_R U
\]
\margin{\begin{tikzcd}[ampersand replacement=\&,row sep=small]
 C\arrow[d,"{T\ \text{``inj''}}"]\\
 \substack{\text{foncteurs contr.}\\ C\to\mathscr G\ \text{qui }\unsure{\text{tr}|\text{ch}}}\arrow[d,"{T'\ \text{``inj''}}"]\\
 \mathscr G^{R^o}\arrow[uu,bend left=65,"T''"]
\end{tikzcd}}
On va montrer que \deleted{si $h$ est un foncteur $C^o\longrightarrow\mathscr G$ du type envisagé, alors} \inserted{$M=T'T(T''M)$} donc $T'T$ est ``surjectif'', donc $T$ et $T'$ sont des ``équivalences''. On doit donc montrer que
\[
 \left\{\begin{aligned}
 &\mathrm{Hom}_C(U,M\otimes_R U)\simeq M\\
 &\text{si}\quad R=\mathrm{Hom}_C(U,U)\\
 &\phantom{\text{si}\quad} M\in\mathscr G^{R^o}
 \end{aligned}\right.
\]
ce qui est laissé au \unsure{soin|soins} du lecteur \ldots

\deleted{\unsure{L}} Ainsi, les foncteurs \deleted{\unsure{covariants}} \inserted{\unsure{contravariant|contravariants}} exacts à gauche $C\longrightarrow\mathscr G$, \deleted{transformant \unsure{une|les}} limites inductives en limites projectives sont essentiellement \inserted{des $h_B$}
% archive sheet number: 15

\scanpage{17}
2. Soient $C,\deleted{\unsure{\widetilde C}}C'$ deux catégories \deleted{additives, abéliennes} \inserted{additives}, $S:\deleted{\unsure{A\longrightarrow\widetilde A}}C\longrightarrow C'$ un foncteur additif de $C$ dans $\deleted{\unsure{\widetilde C}}C'$. On suppose que pour \deleted{tout} \deleted{\unsure{$A\in C$}}, $X\in\bar C$, le foncteur $A\mathrel{\unsure{-|\to|\mapsto}}\mathrm{Hom}_{C'}(S(A),X)$ est un foncteur du type \deleted{$\mathrm{Hom}(A,$} $h_B$ (où $B$ dépend de $X$, mais de façon essentiellement fonctorielle, donc on aura
% slanted
\margin{Mettre en évidence le rôle symétrique joué par $C$ et $C'^{o}$}
\[
 \mathrm{Hom}_{\bar C}(S(\deleted{\unsure{\widetilde A}}A),X)\xleftarrow{\sim}\mathrm{Hom}_C(A,T(X)) \tag*{(1)}
\]
où $T:\deleted{\unsure{\widetilde C}}C'\longrightarrow C$ est un foncteur covariant exact à gauche de $\deleted{\unsure{\widetilde C|C^{o}}}C'$ dans $C$]. Cette condition est satisfaite notamment si $C$ et $\deleted{\unsure{\widetilde C}}C'$ satisfont à $\mathrm{AB}\unsure{(1,|1,|(}3,4,5)$, si $C$ a un générateur, et si de plus \deleted{\unsure{$A\longmapsto\widetilde A$}} $S$ est un foncteur exact à droite (donc $A\mathrel{-}\mathrm{Hom}_C(\deleted{\unsure{\widetilde A}}S(A),X)$ est exact à gauche). \deleted{En particulier,}

En particulier, si $X=S(\deleted{\unsure{\widetilde B}}B)$, on aura
\[
 \mathrm{Hom}_C(A,B)\longrightarrow\mathrm{Hom}_{C'}(\deleted{\unsure{\widetilde A}}S(A),\deleted{\unsure{\widetilde B}}S(B))\xleftarrow{\sim}\mathrm{Hom}_C(A,TS(B)) \tag*{(2)}
\]
Donc faisant $A=B$, on trouve un élément canonique
\[
 \varphi(A)\in\mathrm{Hom}(A,TS(\deleted{\unsure{\widetilde A}}A))\qquad\text{i.e.} \tag*{(3)}
\]
\[
 \varphi:\mathrm{id}_C\longrightarrow TS \tag*{(3 bis)}
\]
D'autre part, on tire de (1), faisant $A=T(Y)$
\[
 \deleted{\mathrm{Hom}_{\widetilde C}(\widetilde A,X)\xrightarrow{\sim}\mathrm{Hom}_C(A,T(X))\longrightarrow\mathrm{Hom}_{\widetilde C}(\widetilde A,\widetilde{T(X)})} \tag*{\deleted{(4)}}
\]
\[
 \mathrm{Hom}_{C'}(\unsure{TS|ST}(\deleted{\unsure{\widetilde Y}}Y),X)\simeq\mathrm{Hom}_C(T(Y),T(X)) \tag*{(4)}
\]
et en particulier, faisant \deleted{\unsure{A}} $X=Y$, on trouve un élément canonique
\[
 \psi(X)\in\mathrm{Hom}_{C'}(ST\deleted{\unsure{\widetilde{\phantom X}}}(X),X)\simeq \tag*{(5)}
\]
\[
 \psi:ST\longrightarrow\deleted{\unsure{\varphi}}\mathrm{id}_{\widetilde C} \tag*{(5 bis)}
\]
% margin note written sideways
\margin{C'est la donnée des $\psi(X)$ qui précise l'homomorphisme (1), puisqu'on a $\mathrm{Hom}_C(A,T(X))\longrightarrow\mathrm{Hom}_{\widetilde C}(\widetilde A,\widetilde{T(X)})\longrightarrow\mathrm{Hom}_{\widetilde C}(A,X)$}
% archive sheet number: 16

\scanpage{18}
\deleted{En plus, on \illegible{1} \unsure{considère} l'\unsure{homom.}}

Considérons \deleted{pour} $A\in C$, d'où \deleted{$\widetilde A\in\widetilde C$ et} $i_A:\deleted{\unsure{A\longrightarrow T\widetilde A}}\inserted{A\mathrel{\unsure{-|\to}}TSA}$ \deleted{\unsure{$\widetilde A$,}} et par suite $S(i_A):SA\longrightarrow STSA$ ; d'autre part, on a aussi $(ST)(SA)\xrightarrow{\unsure{e_{SA}|\varepsilon_{SA}|\psi(SA)}}(SA)$. Ceci dit, on a
\diagram{\begin{tikzcd}[ampersand replacement=\&]
 SA\arrow[r,"{S(\varphi(A))}"]\arrow[rr,bend right=17,"\text{identité}"']
 \& S(TS(A))=(ST)(SA)\arrow[r,"{\psi(SA)}"] \& SA
\end{tikzcd}}
en d'autres termes, $S(A)$ s'identifie par $S(\varphi(A))$ à un facteur direct de $(ST)(S(A))$, la projection naturelle étant $\deleted{e_{S(A)}}\psi(S(A))$. De \textquotedbl on a le diagramme commutatif
\diagram{\begin{tikzcd}[ampersand replacement=\&]
 T(\deleted{\unsure{A}}X) \& T(ST(\deleted{\unsure{A}}X))\deleted{\unsure{\leftarrow}}=\deleted{(TS)}(T(\deleted{\unsure{A}}X))\arrow[l,"{T(\psi(\deleted{\unsure{A}}X))}"'] \&
 T(X)\arrow[l,"{\varphi(T(\unsure{X|A}))}"']\arrow[ll,bend left=13,"\text{identité}"]
\end{tikzcd}}
qui identifie $T(X)$ à un facteur direct de $(TS)(T(X))$.

Pour que l'homomorphisme fonctoriel $\varphi:\mathrm{id}_C\longrightarrow TS$ soit injectif (i.e. les $\varphi(A):A\mathrel{-}TS(A)$ injectifs) il faut et il suffit que les homomorphismes
\[
 S(A,B):\mathrm{Hom}_C(A,B)\longrightarrow\mathrm{Hom}_{C'}(SA,SB)
\]
soient injectifs ; cela signifie aussi que le foncteur $S$ est \deleted{\unsure{fi}} ``injectif''. Cela résulte de la commutativité de
\diagram{\begin{tikzcd}[ampersand replacement=\&]
 \mathrm{Hom}_C(A,B)\arrow[r]\arrow[rr,bend right=12,"{\deleted{\unsure{u}}}"']
 \&\mathrm{Hom}_{C'}(SA,SB) \& \mathrm{Hom}_C(A,TSB)\arrow[l,"\sim"']
\end{tikzcd}}
% archive sheet number: 17
